% GENERATED FILE. Edit canonical lessons, then run npm run generate:books.
% canonical-source-hash: fnv1a64:898a2ff72be90e41
% canonical-lessons: KA-C368-kramena, KA-C368-taruvaya, KA-C368-sarani, KA-C368-ghatta, KA-C368-jarugu, KA-R368-review-362-368

\chapter{Step by Step}
\label{ch:ka-step-by-step}
\hlchaptermodality{368}

\begin{chapteropening}
\textbf{By the end of this chapter:} \emph{I can say gradually, thereafter, a series, a turning point and to take place.}

Complete the last lesson of chapter 368: I can say gradually, thereafter, a series, a turning point and to take place.
\end{chapteropening}

\section[\texorpdfstring{\kn{ಕ್ರಮೇಣ} (kramēṇa)}{ಕ್ರಮೇಣ (kramēṇa)}]{\kn{ಕ್ರಮೇಣ} (kramēṇa) — gradually}

\label{lesson:KA-C368-kramena}



\begin{quote}
Before the new one: say the Kannada for to destroy, then the Kannada for to disappear.
\end{quote}

\subsection*{You'll want to know: \kn{ಕ್ರಮೇಣ}}
\textbf{\kn{ಕ್ರಮೇಣ}} — \emph{kramēṇa} — \textquotedblleft{}gradually\textquotedblright{}.

\begin{quote}
\kn{ಕ್ರಮೇಣ} \kn{ಹಳ್ಳಿ} \kn{ಒಂದು} \kn{ಪಟ್ಟಣವಾಯಿತು}.
\end{quote}

\begin{quote}
\emph{kramēṇa haḷḷi ondu paṭṭaṇavāyitu.}
\end{quote}

\begin{quote}
Gradually the village became a town.
\end{quote}

A formal adverb; in talk people say \kn{ನಿಧಾನವಾಗಿ} or \kn{ಸ್ವಲ್ಪ} \kn{ಸ್ವಲ್ಪವಾಗಿ}.

\begin{cousinweb}
\textbf{\kn{ಕ್ರಮೇಣ}} (\emph{kramēṇa}) is borrowed from Sanskrit \emph{krameṇa}, \textquotedblleft{}in regular course, gradually, by degrees\textquotedblright{} (Monier-Williams), a case form of \emph{krama}, \textquotedblleft{}a step, an order\textquotedblright{}.
\end{cousinweb}

\begin{grammarlens}[title={What you've built}]
One more word for telling what happened.
\end{grammarlens}

\subsection*{Your turn}
\begin{itemize}
\raggedright
  \item \emph{Say it:} \emph{kramēṇa}
  \item \emph{Say it:} \emph{kramēṇa}, once more
  \item \emph{Say it:} \emph{kramēṇa haḷḷi ondu paṭṭaṇavāyitu.} — Gradually the village became a town.
  \item \emph{Recall:} say the Kannada for to destroy, then the Kannada for to disappear, then say \emph{kramēṇa} again
\end{itemize}

\begin{tcolorbox}[breakable,colback=teal!4,colframe=teal!35!black,title={Before you move on}]
What does \kn{ಕ್ರಮೇಣ} mean? (\textquotedblleft{}Gradually\textquotedblright{}.) Say it once more.
\end{tcolorbox}
\section[\texorpdfstring{\kn{ತರುವಾಯ} (taruvāya)}{ತರುವಾಯ (taruvāya)}]{\kn{ತರುವಾಯ} (taruvāya) — thereafter; after}

\label{lesson:KA-C368-taruvaya}



\begin{quote}
Before the new one: say the Kannada for to disappear, then the Kannada for gradually.
\end{quote}

\subsection*{You'll want to know: \kn{ತರುವಾಯ}}
\textbf{\kn{ತರುವಾಯ}} — \emph{taruvāya} — \textquotedblleft{}thereafter; after\textquotedblright{}.

\begin{quote}
\kn{ಯುದ್ಧದ} \kn{ತರುವಾಯ} \kn{ದೇಶ} \kn{ನಿಧಾನವಾಗಿ} \kn{ಬೆಳೆಯಿತು}.
\end{quote}

\begin{quote}
\emph{yuddhada taruvāya dēśa nidhānavāgi beḷeyitu.}
\end{quote}

\begin{quote}
After the war the country grew slowly.
\end{quote}

A written partner of \kn{ನಂತರ}: \kn{ಊಟದ} \kn{ತರುವಾಯ}, after the meal.

\begin{cousinweb}
\textbf{\kn{ತರುವಾಯ}} (\emph{taruvāya}) is native Dravidian: DEDR 3142 lists Kannada \emph{taṟu-vāy} beside Tamil \emph{taṟuvāy}, \textquotedblleft{}an occasion, a particular juncture, a stage\textquotedblright{}.
\end{cousinweb}

\begin{grammarlens}[title={What you've built}]
One more word for telling what happened.
\end{grammarlens}

\subsection*{Your turn}
\begin{itemize}
\raggedright
  \item \emph{Say it:} \emph{taruvāya}
  \item \emph{Say it:} \emph{taruvāya}, once more
  \item \emph{Say it:} \emph{yuddhada taruvāya dēśa nidhānavāgi beḷeyitu.} — After the war the country grew slowly.
  \item \emph{Recall:} say the Kannada for to disappear, then the Kannada for gradually, then say \emph{taruvāya} again
\end{itemize}

\begin{tcolorbox}[breakable,colback=teal!4,colframe=teal!35!black,title={Before you move on}]
What does \kn{ತರುವಾಯ} mean? (\textquotedblleft{}Thereafter; after\textquotedblright{}.) Say it once more.
\end{tcolorbox}
\section[\texorpdfstring{\kn{ಸರಣಿ} (saraṇi)}{ಸರಣಿ (saraṇi)}]{\kn{ಸರಣಿ} (saraṇi) — a series; a string}

\label{lesson:KA-C368-sarani}



\begin{quote}
Before the new one: say the Kannada for gradually, then the Kannada for thereafter.
\end{quote}

\subsection*{You'll want to know: \kn{ಸರಣಿ}}
\textbf{\kn{ಸರಣಿ}} — \emph{saraṇi} — \textquotedblleft{}a series; a string\textquotedblright{}.

\begin{quote}
\kn{ಆ} \kn{ತಿಂಗಳು} \kn{ಊರಿನಲ್ಲಿ} \kn{ಕಳ್ಳತನಗಳ} \kn{ಸರಣಿ} \kn{ನಡೆಯಿತು}.
\end{quote}

\begin{quote}
\emph{ā tiṅgaḷu ūrinalli kaḷḷatanagaḷa saraṇi naḍeyitu.}
\end{quote}

\begin{quote}
That month there was a series of thefts in the town.
\end{quote}

Plural: \kn{ಸರಣಿಗಳು}, \emph{saraṇigaḷu}.

Also a TV or book series, and a cricket series.

\begin{cousinweb}
\textbf{\kn{ಸರಣಿ}} (\emph{saraṇi}) is borrowed from Sanskrit \emph{saraṇi}, \textquotedblleft{}a road, a path, a straight or continuous line\textquotedblright{} (Monier-Williams).
\end{cousinweb}

\begin{grammarlens}[title={What you've built}]
One more word for telling what happened.
\end{grammarlens}

\subsection*{Your turn}
\begin{itemize}
\raggedright
  \item \emph{Say it:} \emph{saraṇi}
  \item \emph{Say it:} \emph{saraṇi}, once more
  \item \emph{Say it:} \emph{ā tiṅgaḷu ūrinalli kaḷḷatanagaḷa saraṇi naḍeyitu.} — That month there was a series of thefts in the town.
  \item \emph{Recall:} say the Kannada for gradually, then the Kannada for thereafter, then say \emph{saraṇi} again
\end{itemize}

\begin{tcolorbox}[breakable,colback=teal!4,colframe=teal!35!black,title={Before you move on}]
What does \kn{ಸರಣಿ} mean? (\textquotedblleft{}A series; a string\textquotedblright{}.) Say it once more.
\end{tcolorbox}
\section[\texorpdfstring{\kn{ಘಟ್ಟ} (ghaṭṭa)}{ಘಟ್ಟ (ghaṭṭa)}]{\kn{ಘಟ್ಟ} (ghaṭṭa) — a turning point; a stage}

\label{lesson:KA-C368-ghatta}



\begin{quote}
Before the new one: say the Kannada for thereafter, then the Kannada for a series.
\end{quote}

\subsection*{You'll want to know: \kn{ಘಟ್ಟ}}
\textbf{\kn{ಘಟ್ಟ}} — \emph{ghaṭṭa} — \textquotedblleft{}a turning point; a stage\textquotedblright{}.

\begin{quote}
\kn{ಆ} \kn{ಪರೀಕ್ಷೆ} \kn{ಅವಳ} \kn{ಜೀವನದ} \kn{ಮುಖ್ಯ} \kn{ಘಟ್ಟ}.
\end{quote}

\begin{quote}
\emph{ā parīkṣe avaḷa jīvanada mukhya ghaṭṭa.}
\end{quote}

\begin{quote}
That exam is a turning point in her life.
\end{quote}

Also a hill slope or a mountain pass: \kn{ಘಟ್ಟದ} \kn{ದಾರಿ}, a road over the hills.

\begin{cousinweb}
\textbf{\kn{ಘಟ್ಟ}} (\emph{ghaṭṭa}) matches Sanskrit \emph{ghaṭṭa}, \textquotedblleft{}a landing place, a quay, a ferry\textquotedblright{} (CDIAL 4414). CDIAL and DEDR 1147 both compare it with Dravidian words for a bank or a dam, such as Kannada \emph{gaṭṭu}, so its first source is not settled; \textquotedblleft{}a slope\textquotedblright{} and \textquotedblleft{}a stage\textquotedblright{} are later senses.
\end{cousinweb}

\begin{grammarlens}[title={What you've built}]
One more word for telling what happened.
\end{grammarlens}

\subsection*{Your turn}
\begin{itemize}
\raggedright
  \item \emph{Say it:} \emph{ghaṭṭa}
  \item \emph{Say it:} \emph{ghaṭṭa}, once more
  \item \emph{Say it:} \emph{ā parīkṣe avaḷa jīvanada mukhya ghaṭṭa.} — That exam is a turning point in her life.
  \item \emph{Recall:} say the Kannada for thereafter, then the Kannada for a series, then say \emph{ghaṭṭa} again
\end{itemize}

\begin{tcolorbox}[breakable,colback=teal!4,colframe=teal!35!black,title={Before you move on}]
What does \kn{ಘಟ್ಟ} mean? (\textquotedblleft{}A turning point; a stage\textquotedblright{}.) Say it once more.
\end{tcolorbox}
\section[\texorpdfstring{\kn{ಜರುಗು} (jarugu)}{ಜರುಗು (jarugu)}]{\kn{ಜರುಗು} (jarugu) — to take place; to shift}

\label{lesson:KA-C368-jarugu}



\begin{quote}
Before the new one: say the Kannada for a series, then the Kannada for a turning point.
\end{quote}

\subsection*{You'll want to know: \kn{ಜರುಗು}}
\textbf{\kn{ಜರುಗು}} — \emph{jarugu} — \textquotedblleft{}to take place; to shift\textquotedblright{}.

\begin{quote}
\kn{ಸಭೆ} \kn{ಶಾಂತವಾಗಿ} \kn{ಜರುಗಿತು}.
\end{quote}

\begin{quote}
\emph{sabhe śāntavāgi jarugitu.}
\end{quote}

\begin{quote}
The meeting took place peacefully.
\end{quote}

Past (it …): \kn{ಜರುಗಿತು}, \emph{jarugitu}.

Events take place: \kn{ಸಮಾರಂಭ} \kn{ಜರುಗಿತು}. It also means to move aside, as \kn{ಸರಿ} does.

\begin{cousinweb}
\textbf{\kn{ಜರುಗು}} (\emph{jarugu}) is native Dravidian: DEDR 2360 lists Kannada \emph{jarugu} beside Tamil \emph{cari}, \textquotedblleft{}to slip away, to slide, to give way\textquotedblright{}.
\end{cousinweb}

\begin{grammarlens}[title={What you've built}]
One more word for telling what happened.
\end{grammarlens}

\subsection*{Your turn}
\begin{itemize}
\raggedright
  \item \emph{Say it:} \emph{jarugu}
  \item \emph{Say it:} \emph{jarugu}, once more
  \item \emph{Say it:} \emph{sabhe śāntavāgi jarugitu.} — The meeting took place peacefully.
  \item \emph{Recall:} say the Kannada for a series, then the Kannada for a turning point, then say \emph{jarugu} again
\end{itemize}

\begin{tcolorbox}[breakable,colback=teal!4,colframe=teal!35!black,title={Before you move on}]
What does \kn{ಜರುಗು} mean? (\textquotedblleft{}To take place; to shift\textquotedblright{}.) Say it once more.
\end{tcolorbox}
\section[(dialogue)]{Looking back: from Builders of Old to Step by Step}

\label{lesson:KA-R368-review-362-368}



\begin{quote}
Say the Kannada for a civilization and to found.
\end{quote}

\subsection*{Your turn}
Say these aloud:

\begin{itemize}
\raggedright
  \item a civilization, to found, to construct, to demolish, a statue
  \item dusk, to rise, to resound, to emerge, to become visible
  \item a diary, an autobiography, a biography, a preface, the gist
  \item a fistfight, violence, a prisoner, to shove, a patrol
  \item an acquittal, a stain, dignity, an oath, to wail
  \item a villain, a hero, an enemy, to destroy, to disappear
  \item gradually, thereafter, a series, a turning point, to take place
\end{itemize}

\begin{tcolorbox}[breakable,colback=teal!4,colframe=teal!35!black,title={Before you move on}]
A civilization? (\textbf{\kn{ನಾಗರಿಕತೆ}}, \emph{nāgarikate}.) To found? (\textbf{\kn{ಸ್ಥಾಪಿಸು}}, \emph{sthāpisu}.) To construct? (\textbf{\kn{ನಿರ್ಮಿಸು}}, \emph{nirmisu}.) To demolish? (\textbf{\kn{ಕೆಡವು}}, \emph{keḍavu}.) A statue? (\textbf{\kn{ಪ್ರತಿಮೆ}}, \emph{pratime}.) Dusk? (\textbf{\kn{ಮುಸ್ಸಂಜೆ}}, \emph{mussañje}.) To rise? (\textbf{\kn{ಮೂಡು}}, \emph{mūḍu}.) To resound? (\textbf{\kn{ಮೊಳಗು}}, \emph{moḷagu}.) To emerge? (\textbf{\kn{ಹೊರಹೊಮ್ಮು}}, \emph{horahommu}.) To become visible? (\textbf{\kn{ಗೋಚರಿಸು}}, \emph{gōcarisu}.) A diary? (\textbf{\kn{ದಿನಚರಿ}}, \emph{dinacari}.) An autobiography? (\textbf{\kn{ಆತ್ಮಕಥೆ}}, \emph{ātmakathe}.) A biography? (\textbf{\kn{ಜೀವನಚರಿತ್ರೆ}}, \emph{jīvanacaritre}.) A preface? (\textbf{\kn{ಪೀಠಿಕೆ}}, \emph{pīṭhike}.) The gist? (\textbf{\kn{ತಿರುಳು}}, \emph{tiruḷu}.) A fistfight? (\textbf{\kn{ಹೊಡೆದಾಟ}}, \emph{hoḍedāṭa}.) Violence? (\textbf{\kn{ಹಿಂಸೆ}}, \emph{hiṃse}.) A prisoner? (\textbf{\kn{ಕೈದಿ}}, \emph{kaidi}.) To shove? (\textbf{\kn{ನೂಕು}}, \emph{nūku}.) A patrol? (\textbf{\kn{ಗಸ್ತು}}, \emph{gastu}.) An acquittal? (\textbf{\kn{ಖುಲಾಸೆ}}, \emph{khulāse}.) A stain? (\textbf{\kn{ಕಳಂಕ}}, \emph{kaḷaṅka}.) Dignity? (\textbf{\kn{ಘನತೆ}}, \emph{ghanate}.) An oath? (\textbf{\kn{ಆಣೆ}}, \emph{āṇe}.) To wail? (\textbf{\kn{ಗೋಳಾಡು}}, \emph{gōḷāḍu}.) A villain? (\textbf{\kn{ಖಳನಾಯಕ}}, \emph{khaḷanāyaka}.) A hero? (\textbf{\kn{ವೀರ}}, \emph{vīra}.) An enemy? (\textbf{\kn{ಶತ್ರು}}, \emph{śatru}.) To destroy? (\textbf{\kn{ನಾಶಮಾಡು}}, \emph{nāśamāḍu}.) To disappear? (\textbf{\kn{ಕಣ್ಮರೆಯಾಗು}}, \emph{kaṇmareyāgu}.) Gradually? (\textbf{\kn{ಕ್ರಮೇಣ}}, \emph{kramēṇa}.) Thereafter? (\textbf{\kn{ತರುವಾಯ}}, \emph{taruvāya}.) A series? (\textbf{\kn{ಸರಣಿ}}, \emph{saraṇi}.) A turning point? (\textbf{\kn{ಘಟ್ಟ}}, \emph{ghaṭṭa}.) To take place? (\textbf{\kn{ಜರುಗು}}, \emph{jarugu}.)
\end{tcolorbox}
